\scanpage{10}

\deleted{\uline{Corollaire}} \uline{Toute déformation infinitésimale \inserted{du premier ordre} de la structure de $C$ qui induit une déformation infinitésimale triviale de la structure $J^0_{C/S}$, est triviale.}

\unsure{Transformant|Transposant|Traduisant} par dualité, cela équivaut à la question si
% slanted
\margin{dégager de façon \unsure{générale} \unsure{ttes} les compatibilités}
\[
H^0(C,\check V\otimes\underline\Omega)=H^0(C,\underline\Omega)\otimes H^0(C,\underline\Omega)\longrightarrow H^0(C,\underline\Omega^{\otimes 2})
\]
(où $\Omega=\Omega^1_{C/k}$) est \uline{surjectif}. Or ce dernier \uline{fait est vrai} \uline{ssi} $C$ \uline{est non hyperelliptique} ; dans le cas \deleted{\unsure{général}} hyperelliptique, on sait que le conoyau de l'hom précédent est de dim 1, donc le conoyau de $H^0(V)\mathbin{\unsure{-|\to}}H^0(N)$ est de dim 1, donc (s'agissant \unsure{du} point $x$ de $M$ \deleted{\unsure{issu}} sur $s\in S$) la fibre de $M$ en $s$ est $\mathrm{Spec}\,k[t]/t^2$.

(ii) On veut montrer que $M\mathbin{\unsure{-|\to}}S$ se factorise en un produit d'un morphisme \uline{propre} par une immersion. Par le critère valuatif, on est ramené à prouver ceci :

\uline{Corollaire} Supposons $S=\mathrm{Spec}(V)$, $V$ anneau de valuation discrète, \deleted{Soit} $\eta$ le pt générique de $S$, \deleted{\unsure{$C_\eta$ une}} On suppose $M\mathbin{\unsure{-|\to}}V$ surjectif i.e. que les fibres de $A$ sur $V$ sont des jacobiennes. \deleted{Alors} Soit $C_\eta$ une courbe sur

% archive sheet number: 9
